\documentclass[a4paper,14pt]{extarticle}

\usepackage{fontspec}
\usepackage{polyglossia}
\setmainlanguage{russian}
\setotherlanguage{english}
\setmainfont[Ligatures=TeX]{Times New Roman}
\newfontfamily\cyrillicfont[Ligatures=TeX]{Times New Roman}
\newfontfamily\cyrillicfonttt{Menlo}[Scale=0.88]
\newfontfamily\cyrillicfontsf{Arial}

\usepackage[
  left=30mm,
  right=15mm,
  top=20mm,
  bottom=20mm
]{geometry}
\usepackage{setspace}
\onehalfspacing
\setlength{\parindent}{1.25cm}
\setlength{\parskip}{0pt}

% Сдаваемая версия ВКР не должна разрывать слова переносами вида
% "обозначе- / ния", особенно на стыке страниц.
\hyphenpenalty=10000
\exhyphenpenalty=10000
\brokenpenalty=10000
\doublehyphendemerits=1000000
\finalhyphendemerits=1000000

\usepackage{indentfirst}
\usepackage{microtype}
\usepackage{graphicx}
\usepackage{adjustbox}
\usepackage{float}
\usepackage{flafter}
\usepackage{caption}
\usepackage{longtable}
\usepackage{booktabs}
\usepackage{array}
\usepackage{calc}
\usepackage{ragged2e}
\usepackage{etoolbox}
\usepackage{enumitem}
\usepackage{amsmath}
\usepackage{amssymb}
\usepackage{fvextra}
\usepackage{listings}
\usepackage{lastpage}
\usepackage{titlesec}
\usepackage{tocloft}
\usepackage{needspace}
\usepackage{placeins}
\usepackage{xurl}
\usepackage{hyperref}
\usepackage{bookmark}
\usepackage{xcolor}

\graphicspath{{./}{../}{docx_pipeline/figures/png/}{docx_pipeline/figures/svg/}}

\hypersetup{
  unicode=true,
  colorlinks=false,
  pdfborder={0 0 0},
  pdftitle={Разработка и внедрение интеллектуальной системы управления микроклиматом в теплице на основе машинного обучения},
  pdfauthor={Ларионов Сергей Иванович}
}

\setcounter{secnumdepth}{0}
\setcounter{tocdepth}{3}

\titleformat{\section}
  {\normalfont\bfseries}
  {}
  {0pt}
  {}
\titlespacing*{\section}{1.25cm}{0pt}{12pt}

\titleformat{\subsection}
  {\normalfont\bfseries}
  {}
  {0pt}
  {}
\titlespacing*{\subsection}{1.25cm}{12pt}{6pt}

\titleformat{\subsubsection}
  {\normalfont\bfseries}
  {}
  {0pt}
  {}
\titlespacing*{\subsubsection}{1.25cm}{10pt}{6pt}

\pretocmd{\section}{\FloatBarrier\clearpage}{}{}
\pretocmd{\subsection}{\Needspace{6\baselineskip}}{}{}
\pretocmd{\subsubsection}{\Needspace{5\baselineskip}}{}{}

\newcommand{\VkrStructuralSection}[1]{%
  \FloatBarrier
  \clearpage
  \phantomsection
  \addcontentsline{toc}{section}{#1}%
  \begin{center}\bfseries #1\end{center}%
  \par\nopagebreak\vspace{12pt}%
}

\newcommand{\VkrStructuralSectionNoToc}[1]{%
  \begin{center}\bfseries #1\end{center}%
  \par\nopagebreak\vspace{12pt}%
}

\newcommand{\VkrAppendixSection}[2]{%
  \FloatBarrier
  \clearpage
  \phantomsection
  \addcontentsline{toc}{section}{ПРИЛОЖЕНИЕ #1 #2}%
  \begin{center}\bfseries ПРИЛОЖЕНИЕ #1\end{center}%
  \par\nopagebreak\vspace{6pt}%
  \begin{center}\bfseries #2\end{center}%
  \par\nopagebreak\vspace{12pt}%
}

\renewcommand{\cftsecleader}{\cftdotfill{\cftdotsep}}
\makeatletter
\renewcommand{\tableofcontents}{%
  \section*{СОДЕРЖАНИЕ}%
  \@starttoc{toc}%
}
\makeatother
\setlength{\cftbeforesecskip}{2pt}
\setlength{\cftbeforesubsecskip}{1pt}

\captionsetup{
  font={normalsize,stretch=1},
  labelformat=empty,
  labelsep=none,
  justification=centering,
  singlelinecheck=false,
  skip=6pt
}
\captionsetup[lstlisting]{
  font={normalsize,stretch=1},
  labelformat=simple,
  labelsep=endash,
  justification=raggedright,
  singlelinecheck=false,
  skip=6pt
}

\setlist[itemize]{label=\textemdash,leftmargin=1.25cm,itemsep=0pt,topsep=0pt,parsep=0pt}
\setlist[enumerate]{leftmargin=1.25cm,itemsep=0pt,topsep=0pt,parsep=0pt}

\lstset{
  basicstyle=\ttfamily\small,
  breaklines=true,
  columns=fullflexible,
  keepspaces=true,
  frame=single,
  framerule=0.4pt,
  framesep=4pt,
  captionpos=t,
  rulecolor=\color{black},
  xleftmargin=0pt,
  xrightmargin=0pt,
  aboveskip=6pt,
  belowskip=6pt
}
\renewcommand{\lstlistingname}{Листинг}
\RecustomVerbatimEnvironment{verbatim}{Verbatim}{
  breaklines=true,
  breakanywhere=true,
  fontsize=\small,
  frame=single,
  framerule=0.4pt,
  framesep=4pt
}
\renewcommand{\texttt}[1]{\textnormal{#1}}

\sloppy
\emergencystretch=2em
\clubpenalty=10000
\widowpenalty=10000
\displaywidowpenalty=10000
\pagestyle{plain}

\providecommand{\tightlist}{%
  \setlength{\itemsep}{0pt}\setlength{\parskip}{0pt}}
\newcommand{\VkrFigureFull}[1]{%
  \adjustbox{max width=\linewidth,max totalheight=0.74\textheight,keepaspectratio,center}{#1}%
}
\newcommand{\VkrFigureInline}[1]{%
  \adjustbox{max width=0.94\linewidth,max totalheight=0.55\textheight,keepaspectratio,center}{#1}%
}
\newcommand{\VkrFigureArchitecture}[1]{%
  \adjustbox{max width=0.98\linewidth,max totalheight=0.74\textheight,keepaspectratio,center}{#1}%
}
\newcommand{\VkrFigureTurnCycle}[1]{%
  \adjustbox{max width=0.98\linewidth,max totalheight=0.76\textheight,keepaspectratio,center}{#1}%
}
\newcommand{\VkrFigureStagePolicy}[1]{%
  \adjustbox{max width=0.98\linewidth,max totalheight=0.65\textheight,keepaspectratio,center}{#1}%
}
\newcommand{\VkrFigureExperimentDesign}[1]{%
  \adjustbox{max width=0.98\linewidth,max totalheight=0.60\textheight,keepaspectratio,center}{#1}%
}
\newcommand{\VkrFigureTall}[1]{%
  \adjustbox{max width=0.88\linewidth,max totalheight=0.50\textheight,keepaspectratio,center}{#1}%
}
\newcommand{\VkrFigureVertical}[1]{%
  \adjustbox{max width=0.88\linewidth,max totalheight=0.56\textheight,keepaspectratio,center}{#1}%
}
\newcommand{\VkrFigureCompact}[1]{%
  \adjustbox{max width=0.92\linewidth,max totalheight=0.48\textheight,keepaspectratio,center}{#1}%
}
\newcommand{\VkrFigureWideCompact}[1]{%
  \adjustbox{max width=0.88\linewidth,max totalheight=0.38\textheight,keepaspectratio,center}{#1}%
}
\providecommand{\pandocbounded}[1]{\VkrFigureFull{#1}}
\providecommand{\passthrough}[1]{#1}
\newcounter{none}

\setlength{\LTpre}{0pt}
\setlength{\LTpost}{6pt}
\setlength{\tabcolsep}{3.5pt}
\setlength{\extrarowheight}{2pt}
\renewcommand{\arraystretch}{1.22}
\AtBeginEnvironment{longtable}{\small\singlespacing}

\newcommand{\VkrTableCaption}[1]{%
  \Needspace{0.44\textheight}%
  \par\smallskip
  {\singlespacing\noindent #1\par}%
  \vspace{2pt}%
}


\begin{document}
\hypersetup{pageanchor=false}
\newcommand{\VkrTitleHeaderLine}[1]{%
  \noindent{\fontsize{12pt}{14pt}\selectfont\makebox[\textwidth][c]{#1}\par}%
}
\newcommand{\VkrUdkField}{УДК \underline{\hspace{1.8cm}}}

\begin{titlepage}
\fontsize{12pt}{14pt}\selectfont
\singlespacing
\VkrTitleHeaderLine{МИНИСТЕРСТВО СЕЛЬСКОГО ХОЗЯЙСТВА РОССИЙСКОЙ ФЕДЕРАЦИИ}

\vspace{0.55cm}

\VkrTitleHeaderLine{ФЕДЕРАЛЬНОЕ ГОСУДАРСТВЕННОЕ БЮДЖЕТНОЕ}
\VkrTitleHeaderLine{ОБРАЗОВАТЕЛЬНОЕ УЧРЕЖДЕНИЕ ВЫСШЕГО ОБРАЗОВАНИЯ}
\VkrTitleHeaderLine{«САНКТ-ПЕТЕРБУРГСКИЙ ГОСУДАРСТВЕННЫЙ}
\VkrTitleHeaderLine{АГРАРНЫЙ УНИВЕРСИТЕТ»}

\vspace{0.75cm}

\VkrTitleHeaderLine{ИНСТИТУТ ЭНЕРГЕТИКИ}

\vspace{0.75cm}

\VkrTitleHeaderLine{КАФЕДРА ТЕПЛОТЕХНИКИ И ЭНЕРГООБЕСПЕЧЕНИЯ}

\vspace{0.9cm}

\begin{flushright}
На правах рукописи

\vspace{0.2cm}

\VkrUdkField
\end{flushright}

\vspace{0.75cm}

\begin{center}
Ларионов Сергей Иванович

\vspace{0.8cm}

\textbf{«Разработка и внедрение интеллектуальной системы управления микроклиматом в теплице на основе машинного обучения для повышения эффективности выращивания растений»}

\vspace{0.85cm}

Выпускная квалификационная работа магистра

\vspace{0.65cm}

Направление подготовки\\
«Теплотехника и теплоэнергетика в агропромышленном комплексе»
\end{center}

\vspace{0.4cm}

\begin{flushright}
\begin{minipage}{0.46\textwidth}
Выпускная квалификационная\\
работа защищена

\vspace{0.15cm}

«\underline{\hspace{0.8cm}}»\underline{\hspace{2.5cm}}2026 г.

\vspace{0.18cm}

Оценка

\vspace{0.12cm}

\underline{\hspace{3.0cm}}

\vspace{0.18cm}

Секретарь ГЭК \underline{\hspace{2.2cm}}
\end{minipage}
\end{flushright}

\vfill

\begin{center}
г. Санкт-Петербург\\
2026
\end{center}
\end{titlepage}

\setcounter{page}{2}
\hypersetup{pageanchor=true}

\VkrStructuralSectionNoToc{РЕФЕРАТ}\label{ux440ux435ux444ux435ux440ux430ux442}

Работа содержит \pageref{LastPage} с., 0 рис., 1 табл., 32 источн.

Выпускная квалификационная работа посвящена разработке интеллектуальной
системы управления микроклиматом теплицы на основе машинного обучения.
Теплица рассматривается как теплотехнический и киберфизический объект, в
котором температурный режим, влажность, воздухообмен и концентрация
углекислого газа должны анализироваться совместно с управляющими
воздействиями.

Цель работы состоит в обосновании и разработке расширяемой
IoT/ML-платформы для мониторинга, журналирования и последующего
прогнозного управления микроклиматом теплицы. В текущей редакции
основное внимание уделено постановке задачи, анализу литературы и
обоснованию архитектурной рамки, в которой разнородные датчики и
исполнительные устройства подключаются через единую модель обмена
данными.

В аналитической части рассмотрены современные подходы к smart
greenhouse, IoT-архитектурам, прогнозированию микроклимата,
предиктивному управлению, обучению с подкреплением, энергоэффективности
и обработке неоднородных потоков сенсорных данных. На основе обзора
сформулированы требования к системе как к расширяемой платформе,
способной не только фиксировать параметры микроклимата, но и накапливать
данные для дальнейшего применения методов машинного обучения.

Ключевые слова: теплица, микроклимат, тепловлажностный режим,
интеллектуальное управление, интернет вещей, машинное обучение,
прогнозирование, энергоэффективность.

\clearpage

\VkrStructuralSectionNoToc{СОДЕРЖАНИЕ}
\makeatletter\@starttoc{toc}\makeatother
\clearpage

\VkrStructuralSection{ПЕРЕЧЕНЬ СОКРАЩЕНИЙ И ОБОЗНАЧЕНИЙ}\label{ux441ux43fux438ux441ux43eux43a-ux441ux43eux43aux440ux430ux449ux435ux43dux438ux439-ux438-ux43eux431ux43eux437ux43dux430ux447ux435ux43dux438ux439}

{\def\LTcaptype{none} % do not increment counter
\begin{longtable}[]{@{}
  >{\RaggedRight\arraybackslash}p{(\linewidth - 4\tabcolsep) * \real{0.3333}}
  >{\RaggedRight\arraybackslash}p{(\linewidth - 4\tabcolsep) * \real{0.3333}}
  >{\RaggedRight\arraybackslash}p{(\linewidth - 4\tabcolsep) * \real{0.3333}}@{}}
\hline
\begin{minipage}[b]{\linewidth}\RaggedRight
Сокращение
\end{minipage} & \begin{minipage}[b]{\linewidth}\RaggedRight
Расшифровка
\end{minipage} & \begin{minipage}[b]{\linewidth}\RaggedRight
Пояснение в контексте работы
\end{minipage} \\
\hline
\endhead
\hline
\endlastfoot
АСУ & автоматизированная система управления & система наблюдения и
управления параметрами микроклимата \\
ВКР & выпускная квалификационная работа & итоговая квалификационная
работа магистранта \\
ИИ & искусственный интеллект & методы анализа данных и поддержки
принятия решений \\
IoT & internet of things & подход к объединению датчиков, исполнительных
устройств и программных сервисов \\
ML & machine learning & машинное обучение, применяемое для анализа и
прогнозирования временных рядов \\
MQTT & message queuing telemetry transport & легкий протокол обмена
сообщениями в IoT-системах \\
CO\(_2\) & carbon dioxide & концентрация углекислого газа в воздухе
теплицы \\
MPC & model predictive control & модельно-предиктивное управление \\
RL & reinforcement learning & обучение с подкреплением \\
MAE & mean absolute error & средняя абсолютная ошибка прогноза \\
RMSE & root mean squared error & корень из средней квадратичной ошибки
прогноза \\
\end{longtable}
}

\VkrStructuralSection{ВВЕДЕНИЕ}\label{ux432ux432ux435ux434ux435ux43dux438ux435}

Современные тепличные комплексы требуют точного контроля параметров
микроклимата: температуры, влажности, концентрации CO\(_2\), режимов
вентиляции, нагрева, увлажнения и полива. С теплотехнической точки
зрения теплица является объектом с тепловым балансом, влажностным
режимом и вентиляционным воздухообменом. Даже в малой теплице параметры
изменяются не изолированно: включение вентилятора влияет на температуру
и влажность, открытие заслонок меняет воздухообмен, работа нагревателей
имеет инерционный эффект, а концентрация CO\(_2\) зависит от вентиляции,
фотосинтетической активности растений и притока наружного воздуха.
Поэтому управление микроклиматом нельзя сводить только к ручному
включению оборудования или простым пороговым действиям без накопления
истории.

Актуальность работы связана с переходом тепличных систем от локальной
автоматики к киберфизическим IoT-решениям. В современных исследованиях
smart greenhouse рассматривается как система, объединяющая датчики,
исполнительные механизмы, сетевой обмен, хранилище данных, веб-интерфейс
и интеллектуальный аналитический слой {[}1; 2{]}. Такая постановка
особенно важна для малых и учебно-экспериментальных теплиц: если система
проектируется только под один датчик и одно реле, она быстро перестает
быть исследовательской платформой. Для дальнейшего применения машинного
обучения требуется расширяемая архитектура, в которой новые сенсорные
узлы, исполнительные каналы и расчетные признаки подключаются без
переписывания всей системы.

Отдельное значение имеет энергоэффективность. Нагрев, вентиляция,
увлажнение и освещение потребляют значительные ресурсы, а вентиляционные
воздействия могут приводить к теплопотерям {[}28; 31{]}. При этом
реакция теплицы на управляющее воздействие зависит от текущего состояния
среды, внешних условий, инерции конструкций, плотности растений и режима
воздухообмена. Поэтому управление должно учитывать не только текущие
значения датчиков, но и прогноз развития микроклимата {[}14; 25{]}.

Методы машинного обучения позволяют использовать накопленную телеметрию
для прогнозирования температуры, влажности, CO\(_2\), роста растений и
урожайности. В литературе применяются LSTM, RNN, TCN, Transformer и
вероятностные модели, ориентированные на временные ряды тепличных данных
{[}9; 10; 11; 13{]}. Однако для практического применения таких моделей
сначала необходим надежный инженерный слой: датчики должны стабильно
публиковать данные, исполнительные устройства должны управляться через
понятный интерфейс, а история измерений и воздействий должна сохраняться
в пригодном для анализа виде.

В рамках данной работы рассматривается разработка и опытное внедрение
расширяемой программно-аппаратной IoT/ML-платформы мониторинга и
управления микроклиматом теплицы. Платформа объединяет микроконтроллеры,
MQTT-брокер, веб-панель оператора, правила автоматизации, профили
управления, журналирование данных и подготовку истории для машинного
обучения. Реальная история измерений и управляющих воздействий
необходима для перехода от эмпирического управления к имитационной или
ML-модели тепловлажностного режима {[}32; 29{]}. Машинное обучение в
работе рассматривается как верхний интеллектуальный слой, который
опирается на данные, собранные инженерной системой. Такой подход
позволяет сохранить надежность базового управления и одновременно
подготовить основу для прогнозных моделей и дальнейшего
интеллектуального управления.

Цель работы: разработать и внедрить расширяемую IoT/ML-платформу
мониторинга и управления тепловлажностным режимом теплицы с возможностью
подключения разнородных датчиков, управления исполнительными
устройствами, накопления истории и последующего применения методов
машинного обучения для прогнозирования и интеллектуальной поддержки
управления.

Для достижения цели необходимо решить следующие задачи:

\begin{enumerate}
\def\labelenumi{\arabic{enumi})}
\tightlist
\item
  Проанализировать современные подходы к автоматизации микроклимата
  теплиц, IoT-архитектурам, прогнозированию микроклимата и
  интеллектуальному управлению.
\item
  Определить требования к расширяемой системе мониторинга и управления:
  состав параметров, датчики, исполнительные механизмы, канал связи,
  хранение данных, интерфейс оператора и пригодность истории для
  ML-анализа.
\item
  Спроектировать архитектуру системы на базе микроконтроллеров, MQTT,
  серверного приложения, базы данных истории и единой модели обмена
  сообщениями.
\item
  Реализовать программно-аппаратный контур: сенсорные узлы, релейный
  контроллер, MQTT-обмен, dashboard, журналирование показаний и
  состояний.
\item
  Подготовить основу для интеллектуального слоя: накопление временных
  рядов, экспорт данных, правила и профили управления, план
  ML-эксперимента.
\item
  Проверить работоспособность системы на опытном объекте, сформировать
  датасет эксплуатации и определить направления дальнейшего развития.
\end{enumerate}

Объект исследования: тепловлажностный режим теплицы как управляемая
киберфизическая система агропромышленного комплекса.

Предмет исследования: методы и программно-аппаратные средства
мониторинга, управления и интеллектуальной поддержки регулирования
микроклимата теплицы.

Практическая значимость работы состоит в создании системы, которая может
использоваться как основа для эксплуатации малой теплицы и для
дальнейших экспериментов с прогнозированием микроклимата. Реализованная
архитектура позволяет собирать телеметрию, управлять исполнительными
устройствами, хранить историю и формировать датасет для машинного
обучения.

Научная новизна работы заключается в архитектурно-методическом подходе к
построению малой теплицы как расширяемой IoT/ML-платформы. В отличие от
одноразовой схемы «датчик - реле - панель управления», в работе
рассматривается единый контур, где разнородные датчики публикуют данные
в общей модели обмена, управляющие воздействия регистрируются совместно
с параметрами микроклимата, а накопленная история становится основой для
последующего ML-прогнозирования реакции теплицы. При этом не заявляется
разработка нового алгоритма машинного обучения: интеллектуальный слой
рассматривается как прогнозная надстройка над безопасным правиловым
контуром.

Структура работы включает введение, три главы, заключение, список
литературы и приложения. В первой главе приведен аналитический обзор
источников по smart greenhouse, IoT, прогнозированию и методам
управления. Во второй главе описано проектирование системы. В третьей
главе рассмотрены реализация, проверка работоспособности и подготовка
ML-эксперимента.

\section{1 Аналитический обзор систем интеллектуального управления
микроклиматом
теплицы}\label{ux430ux43dux430ux43bux438ux442ux438ux447ux435ux441ux43aux438ux439-ux43eux431ux437ux43eux440-ux441ux438ux441ux442ux435ux43c-ux438ux43dux442ux435ux43bux43bux435ux43aux442ux443ux430ux43bux44cux43dux43eux433ux43e-ux443ux43fux440ux430ux432ux43bux435ux43dux438ux44f-ux43cux438ux43aux440ux43eux43aux43bux438ux43cux430ux442ux43eux43c-ux442ux435ux43fux43bux438ux446ux44b}

\subsection{1.1 Теплица как киберфизическая система
управления}\label{ux442ux435ux43fux43bux438ux446ux430-ux43aux430ux43a-ux43aux438ux431ux435ux440ux444ux438ux437ux438ux447ux435ux441ux43aux430ux44f-ux441ux438ux441ux442ux435ux43cux430-ux443ux43fux440ux430ux432ux43bux435ux43dux438ux44f}

Современная теплица является не просто защищенным сооружением для
выращивания растений, а киберфизической системой, в которой датчики,
исполнительные механизмы, микроконтроллеры, серверное программное
обеспечение и алгоритмы принятия решений работают как единый контур.
Целевыми параметрами такого контура выступают температура воздуха,
относительная влажность, концентрация CO\(_2\), освещенность, влажность
почвы или субстрата, режим полива и состояние растений. Поддержание этих
параметров влияет на урожайность, скорость роста, устойчивость растений
к стрессу и энергоемкость выращивания.

В обзорной литературе последних лет тепличное хозяйство все чаще
рассматривается через призму smart greenhouse, то есть управляемой
среды, где IoT, искусственный интеллект, спектральные датчики,
мультимодальное слияние данных и интеллектуальные системы используются
для повышения эффективности выращивания. В обзоре El Ouaham et
al.~показано, что развитие smart greenhouse связано не только с заменой
ручного контроля на датчики, но и с переходом к системам, которые
собирают разнородные данные, анализируют их и поддерживают принятие
решений {[}1{]}. Для данной ВКР это принципиально: ценность системы
определяется не только тем, что она измеряет температуру и включает
реле, а тем, что она создает расширяемый контур данных для последующего
прогнозирования и анализа реакции теплицы.

Русскоязычные источники подтверждают ту же постановку на инженерном
уровне. В сравнительном анализе современных систем автоматизированного
управления микроклиматом подчеркивается, что автоматизация теплиц
направлена на снижение эксплуатационных затрат, повышение урожайности и
переход от простого поддержания параметров к интеллектуальному
управлению {[}3{]}. Халина и Цыбанюк описывают базовую структуру
автоматизированной системы: датчики измеряют параметры среды, блок
управления сравнивает их с заданными значениями, после чего
исполнительные механизмы изменяют температуру, влажность, освещение или
другие параметры {[}4{]}. Такая схема является исходной, но в
современных условиях она требует расширения прогнозными и аналитическими
функциями.

Особенность тепличного микроклимата состоит в инерционности и
взаимосвязанности параметров. Нагрев, вентиляция, увлажнение, подача
CO\(_2\) и освещение влияют друг на друга, а реакция растений
проявляется не мгновенно. Тарков и Сукьясов рассматривают задачу
управления параметрами микроклимата через уравнения баланса и отмечают
проблему запаздывания процесса регулирования {[}5{]}. Это означает, что
реактивное управление по текущему отклонению от уставки не всегда
обеспечивает оптимальный результат: система может опоздать с
воздействием или создать избыточные энергетические затраты.

\subsection{1.2 Тепловлажностный режим и энергетическая рамка
теплицы}\label{ux442ux435ux43fux43bux43eux432ux43bux430ux436ux43dux43eux441ux442ux43dux44bux439-ux440ux435ux436ux438ux43c-ux438-ux44dux43dux435ux440ux433ux435ux442ux438ux447ux435ux441ux43aux430ux44f-ux440ux430ux43cux43aux430-ux442ux435ux43fux43bux438ux446ux44b}

Для направления «Теплотехника и теплоэнергетика в агропромышленном
комплексе» принципиально важно рассматривать теплицу как объект тепло- и
массообмена. Температура, влажность и CO\(_2\) формируются под действием
солнечной радиации, теплопотерь через ограждения, вентиляционного
воздухообмена, испарения и транспирации растений, нагрева, увлажнения и
работы вспомогательного оборудования. В обзоре Chen et al.~тепличная
система прямо описывается как сложная многовходовая и многовыходовая
система, где управление температурой, влажностью, светом, CO\(_2\) и
воздушными потоками должно одновременно обеспечивать рост растений и
энергосбережение {[}28{]}.

Динамические модели микроклимата теплицы обычно опираются на баланс
энергии, энтальпии и массы. Это важно для текущей ВКР: даже если
практическая реализация использует датчики и реле, физическая сущность
задачи остается теплотехнической. Самарин и др. формулируют задачу
имитационного моделирования оптимального управления микроклиматом
сельскохозяйственного объекта через температурно-влажностный режим,
тепловой баланс, влажностное состояние конструкций и энергосберегающий
режим выращивания {[}32{]}. Такой подход поддерживает выбранную рамку:
система должна не только включать устройства, но и формировать данные
для модели поведения теплицы.

Отдельная проблема состоит в пространственной неоднородности
микроклимата. Ogunlowo et al.~показывают, что распределение тепла и
массы в одно- и многопролетных теплицах влияет на точность оценки
энергетической потребности; при игнорировании распределения параметров
возможна переоценка или недооценка расчетов {[}29{]}. Это обосновывает
применение нескольких измерительных узлов внутри теплицы и последующий
анализ расхождений между датчиками, а не только использование одного
среднего значения.

Вентиляция является одним из ключевых каналов воздействия на
микроклимат. Она одновременно снижает температуру, удаляет влагу,
изменяет воздухообмен и влияет на концентрацию CO\(_2\). Li et al.~на
данных девяти односкатных теплиц показывают, что открытие вентиляции
меняет температуру, относительную влажность, VPD и CO\(_2\), а также
зависит от плотности посадки, высоты растений, внутренних покрытий и
времени открытия вентиляционного проема {[}30{]}. Поэтому события
включения вентиляторов и приточных каналов в собственной системе должны
журналироваться вместе с температурой, влажностью и CO\(_2\): без этого
невозможно корректно анализировать реакцию теплицы на управляющее
воздействие.

Энергетический аспект особенно заметен при управлении вентиляцией и
отоплением. Ghaderi et al.~отмечают, что вентиляционные потери тепла
являются одним из значимых факторов энергетической эффективности теплиц,
а климат-контроль может составлять основную долю энергопотребления
тепличного хозяйства {[}31{]}. Следовательно, интеллектуальный слой
управления должен стремиться не только удерживать параметры в допустимом
диапазоне, но и снижать лишние переключения, преждевременные теплопотери
и компенсирующие действия после запаздывания.

\subsection{1.3 IoT-архитектура и расширяемость тепличных
систем}\label{iot-ux430ux440ux445ux438ux442ux435ux43aux442ux443ux440ux430-ux438-ux440ux430ux441ux448ux438ux440ux44fux435ux43cux43eux441ux442ux44c-ux442ux435ux43fux43bux438ux447ux43dux44bux445-ux441ux438ux441ux442ux435ux43c}

Практическое внедрение интеллектуального управления начинается с
надежного сбора телеметрии. IoT-подходы позволяют связать датчики,
исполнительные устройства, контроллеры и программные сервисы в единую
инфраструктуру. Bersani et al.~рассматривают smart greenhouse в
контексте Industry 4.0 и показывают, что IoT-системы для теплиц включают
сенсорные сети, каналы связи, удаленный мониторинг и управляющие контуры
{[}2{]}. Такой подход хорошо согласуется с текущим проектом, где уже
присутствуют прошивки микроконтроллеров и серверный dashboard.

В инженерных реализациях часто используются микроконтроллеры, Raspberry
Pi, Arduino/ESP, MQTT, веб-интерфейсы и базы данных. Кельсин в ВКР по
автоматизированной системе управления тепличным хозяйством рассматривает
архитектуру на микроконтроллерах, близкую к практическим учебным и малым
производственным системам {[}7{]}. Статья по MQTT-based smart greenhouse
monitoring дополнительно показывает, почему MQTT удобен для телеметрии:
это легкий протокол обмена сообщениями, пригодный для IoT-устройств с
ограниченными ресурсами {[}23{]}. Для текущей ВКР это важно, поскольку
коммуникационный слой должен быть достаточно простым для
микроконтроллеров и достаточно надежным для накопления данных.

Для перехода от локальной автоматики к платформе недостаточно просто
добавить веб-интерфейс. Система должна иметь расширяемую модель
подключаемых устройств: разные датчики могут измерять температуру,
влажность, CO\(_2\), освещенность, влажность субстрата или
дополнительные расчетные параметры, но для верхнего уровня они должны
попадать в историю в сопоставимой структуре. Такой подход согласуется с
направлением multimodal data fusion в smart greenhouse: данные от разных
источников становятся не разрозненными показаниями, а общей
информационной базой для анализа состояния растений и микроклимата
{[}1{]}. В практической архитектуре это означает необходимость
идентификатора устройства, типа метрики, времени измерения, единиц
измерения и признака качества данных.

Расширяемость важна и для исполнительных устройств. В малой теплице реле
может управлять вентилятором, заслонкой, насосом, увлажнителем или
нагревателем, но для анализа микроклимата все эти действия должны
фиксироваться как события воздействия на тепловлажностный режим. Без
совместной регистрации датчиков и реле невозможно отличить естественное
изменение температуры от реакции на вентиляцию или нагрев. Поэтому
журнал исполнительных воздействий является не вспомогательной функцией
dashboard, а обязательной частью ML-ready датасета.

Прикладные smart farming эксперименты демонстрируют, что тепличная
система должна рассматриваться как полный стек: датчики, исполнительные
механизмы, логика управления, хранение данных и интерфейс оператора.
Joni et al.~описывают IoT-enhanced greenhouse design для выращивания
риса, где IoT-инфраструктура связывает параметры среды и агротехническую
задачу {[}24{]}. Несмотря на отличие культуры, работа полезна как пример
построения комплексной системы. Thwin et al.~идут еще дальше и связывают
прогноз микроклимата с web integration и adaptive equipment control
{[}16{]}. Это особенно важно для данной ВКР, так как собственный проект
уже содержит dashboard, прошивки и историю событий, а значит может быть
описан как первый инкремент интеллектуального контура, а не как
завершенная одноцелевая автоматика.

Отдельно следует отметить работы, где IoT-инфраструктура связывается с
reinforcement learning. Platero-Horcajadas et al.~рассматривают
интеграцию IoT и RL для оптимизированного climate control {[}21{]}. Для
текущей ВКР этот источник важен не как требование немедленно реализовать
RL-агента, а как подтверждение архитектурного тезиса: интеллектуальный
регулятор возможен только там, где есть непрерывная телеметрия,
наблюдаемые исполнительные воздействия и воспроизводимый обмен данными.

\subsection{1.4 Прогнозирование микроклимата, роста и
урожайности}\label{ux43fux440ux43eux433ux43dux43eux437ux438ux440ux43eux432ux430ux43dux438ux435-ux43cux438ux43aux440ux43eux43aux43bux438ux43cux430ux442ux430-ux440ux43eux441ux442ux430-ux438-ux443ux440ux43eux436ux430ux439ux43dux43eux441ux442ux438}

Накопление телеметрии само по себе не делает систему интеллектуальной.
Интеллектуальный слой появляется тогда, когда исторические данные
используются для прогнозирования будущего состояния микроклимата, роста
растений или урожайности. Alhnaity et al.~применяют LSTM для
прогнозирования роста и урожайности в greenhouse environments и
показывают, что последовательностные модели способны извлекать полезные
зависимости из временных рядов тепличных данных {[}9{]}. Gong et
al.~предлагают связку TCN и RNN для прогнозирования урожайности и
сравнивают ее с классическими ML-подходами {[}10{]}. Эти работы
обосновывают использование истории наблюдений, а не только текущего
значения датчика.

Для задачи управления микроклиматом особенно важны работы, где
прогнозируются непосредственно температура, влажность и CO\(_2\). Ahn et
al.~сравнивают transformer-based и RNN-based модели для предсказания
температуры, относительной влажности и концентрации CO\(_2\) в теплице
{[}11{]}. Huang et al.~используют Attention CNN-LSTM для прогнозирования
среды грибной теплицы, также работая с температурой, влажностью и
CO\(_2\) {[}12{]}. Эти источники показывают, что микроклимат является
естественным объектом для time-series prediction, а не только для
порогового контроля. При этом качество прогноза зависит не только от
выбора модели, но и от состава признаков: полезными могут быть лаги
датчиков, время суток, сезонность, состояние вентиляции, нагрева и
увлажнения, а также различия между точками измерения.

Отдельное направление связано с прогнозным управлением и
энергоэффективностью. Aborujilah et al.~рассматривают forecast-driven
climate control для smart greenhouse и используют LSTM-модель для
proactive adjustments, направленных на снижение энергетических затрат
{[}14{]}. Soumo et al.~описывают data-driven greenhouse climate
regulation при выращивании салата с использованием BiLSTM и GRU
predictive control {[}15{]}. Thwin et al.~предлагают multivariate
multistep LSTM, который прогнозирует микроклимат и взаимодействует с
системой управления оборудованием через веб-интеграцию {[}16{]}. В
совокупности эти работы позволяют обосновать, что прогнозный слой может
быть включен в практическую АСУ теплицы.

При этом прогнозы не следует рассматривать как абсолютно надежные. Choi
and Yang предлагают probabilistic deep learning framework для greenhouse
microclimate prediction с учетом time-varying uncertainty и covariance
analysis {[}13{]}. Для диплома этот источник важен методологически: даже
если в прототипе применяется более простая модель, в тексте необходимо
признавать неопределенность прогноза и необходимость защитных
ограничений нижнего уровня управления.

\subsection{1.5 Методы управления: PID, fuzzy, MPC, RL и эвристическая
оптимизация}\label{ux43cux435ux442ux43eux434ux44b-ux443ux43fux440ux430ux432ux43bux435ux43dux438ux44f-pid-fuzzy-mpc-rl-ux438-ux44dux432ux440ux438ux441ux442ux438ux447ux435ux441ux43aux430ux44f-ux43eux43fux442ux438ux43cux438ux437ux430ux446ux438ux44f}

Классические системы управления теплицей часто строятся на пороговых
правилах, PID-регуляторах или нечеткой логике. Их преимущество состоит в
простоте, интерпретируемости и возможности реализации на
микроконтроллерах. Воробьева и др. рассматривают интеллектуальную
систему контроля температуры в теплице и сравнивают PID-регулирование с
нечетким регулятором {[}6{]}. Для инженерного проекта такие методы
остаются важными, потому что нижний уровень управления должен быть
предсказуемым и безопасным.

Однако при усложнении задачи возникает потребность в методах, которые
учитывают будущую динамику, ограничения и стоимость ресурсов. Model
Predictive Control является одним из основных подходов в научной
литературе по greenhouse climate control, поскольку он использует модель
системы для оптимизации управляющих воздействий на горизонте прогноза.
Mallick et al.~показывают, что MPC явно учитывает физические
ограничения, но сильно зависит от точности модели; для компенсации
неопределенности они предлагают связку MPC и reinforcement learning
{[}17{]}. Msaad et al.~также рассматривают RL-guided MPC для автономного
управления теплицей {[}18{]}.

Сравнение MPC и RL важно для корректного позиционирования ВКР. Morcego
et al.~рассматривают MPC как model-based подход, а RL как learning-based
подход к greenhouse climate control {[}19{]}. MPC лучше интерпретируется
и естественно работает с ограничениями, тогда как RL способен улучшать
стратегию на основе данных, но требует осторожности из-за рисков
обучения и безопасности. Поэтому для практической ВКР разумно описывать
RL/MPC не как обязательную реализацию в полном объеме, а как верхний
интеллектуальный слой, который может выдавать рекомендации или
корректировать уставки, тогда как нижний контур остается защищенным.

Перспективным направлением является включение оператора или агронома в
цикл обучения. Xiao et al.~предлагают grower-in-the-loop interactive
reinforcement learning для greenhouse climate control {[}20{]}. Такой
подход хорошо соответствует практической логике внедрения: система может
не заменять оператора полностью, а поддерживать его решения и
использовать экспертную обратную связь.

Помимо MPC/RL, в литературе встречаются эвристические оптимизационные
методы. Jawad et al.~используют artificial bee colony algorithm для
energy optimization and plant comfort management в smart greenhouse
{[}22{]}. Этот источник показывает, что задача управления микроклиматом
может решаться не только нейросетями или MPC, но и другими
оптимизационными методами. Для текущей работы это полезно как обзор
альтернатив, но основная логика ВКР остается связанной с IoT,
прогнозированием и гибридным управлением.

\subsection{1.6 Энергоэффективность и устойчивость тепличного
управления}\label{ux44dux43dux435ux440ux433ux43eux44dux444ux444ux435ux43aux442ux438ux432ux43dux43eux441ux442ux44c-ux438-ux443ux441ux442ux43eux439ux447ux438ux432ux43eux441ux442ux44c-ux442ux435ux43fux43bux438ux447ux43dux43eux433ux43e-ux443ux43fux440ux430ux432ux43bux435ux43dux438ux44f}

Энергопотребление является одним из ключевых ограничений тепличного
производства. Отопление, вентиляция, увлажнение, освещение и подача
CO\(_2\) требуют ресурсов, а реактивное управление может приводить к
избыточным включениям исполнительных механизмов. Aborujilah et al.~прямо
указывают на проблему реактивных стратегий и предлагают forecast-driven
подход с LSTM для снижения энергетических потерь {[}14{]}. Эту же линию
поддерживает теплотехническая литература по вентиляционным потерям и
интеграции отопления, охлаждения и вентиляции {[}31{]}.

Более широкий энергетический контекст раскрывается в обзоре Wu and Fu по
Agricultural Energy Internet {[}25{]}. В нем greenhouse environmental
control рассматривается как энергетически управляемый процесс наряду с
ирригацией, дополнительным освещением, логистикой и
сельскохозяйственными машинами. Для ВКР этот источник полезен в разделе
актуальности: интеллектуальное управление микроклиматом не только
улучшает условия выращивания, но и связано с задачей рационального
использования энергии.

Энергоэффективность также связана с архитектурой управления. Если
система умеет прогнозировать состояние микроклимата, она может заранее
выбирать более мягкие воздействия, снижать частоту резких переключений и
избегать компенсирующих действий после запаздывания. Однако такие
преимущества возможны только при наличии надежной телеметрии,
устойчивого обмена данными и защитных правил нижнего уровня.

\subsection{1.7 Выводы по аналитическому
обзору}\label{ux432ux44bux432ux43eux434ux44b-ux43fux43e-ux430ux43dux430ux43bux438ux442ux438ux447ux435ux441ux43aux43eux43cux443-ux43eux431ux437ux43eux440ux443}

Проведенный обзор показывает, что тема интеллектуального управления
микроклиматом теплицы находится на пересечении нескольких направлений:
теплотехники, тепло- и массообмена, IoT-архитектур, автоматизированного
управления, машинного обучения, прогнозирования временных рядов,
энергоэффективности и человеко-ориентированного принятия решений.
Практические инженерные источники описывают датчики, контроллеры, MQTT,
dashboard и исполнительные механизмы, но часто ограничиваются пороговыми
или простыми регуляторами. Современные научные работы предлагают LSTM,
RNN, Transformer, Attention CNN-LSTM, вероятностные модели, MPC и RL,
однако не всегда доводят эти методы до конкретной малой теплицы с
реальной прошивкой, пользовательским интерфейсом и совместным
журналированием измерений и воздействий.

На этом фоне текущая ВКР может быть обоснована как разработка и опытное
внедрение расширяемой IoT/ML-платформы управления микроклиматом, которая
создает основу для интеллектуального слоя. Нижний уровень системы должен
обеспечивать сбор данных, безопасное управление и связь с
исполнительными устройствами. Верхний уровень должен накапливать
телеметрию, анализировать временные ряды, формировать прогнозы и
поддерживать корректировку уставок. Научный разрыв, закрываемый работой,
состоит не в создании нового ML-алгоритма, а в переходе от разрозненной
автоматики к воспроизводимой платформе: разнородные датчики подключаются
через единую модель обмена, события реле регистрируются вместе с
микроклиматом, а история становится пригодной для ML-прогнозирования
реакции теплицы. Такой гибридный подход согласует надежность
практической АСУ с современными направлениями ML/MPC/RL и позволяет
постепенно развивать систему без риска нереалистичного обещания
полностью автономного AI-управления.

\VkrStructuralSection{СПИСОК ИСПОЛЬЗОВАННЫХ ИСТОЧНИКОВ}\label{ux441ux43fux438ux441ux43eux43a-ux438ux441ux43fux43eux43bux44cux437ux43eux432ux430ux43dux43dux44bux445-ux438ux441ux442ux43eux447ux43dux438ux43aux43eux432}

\begin{enumerate}
\def\labelenumi{\arabic{enumi}.}
\tightlist
\item
  El Ouaham W., Sadik M., Ennajih A., Mouzouna Y., Orchi H., Elouaham S.
  Smart Greenhouses in the Era of IoT and AI: A Comprehensive Review of
  AI Applications, Spectral Sensing, Multimodal Data Fusion, and
  Intelligent Systems // Agriculture. 2026. Vol. 16, No.~7. Article 761.
  DOI: 10.3390/agriculture16070761.
\item
  Bersani C., Ruggiero C., Sacile R., Soussi A., Zero E. Internet of
  Things Approaches for Monitoring and Control of Smart Greenhouses in
  Industry 4.0 // Energies. 2022. Vol. 15, No.~10. Article 3834. DOI:
  10.3390/en15103834.
\item
  Ван Вэйянь, Харисов А. Р. Сравнительный анализ современных систем
  автоматизированного управления микроклиматом в теплице // Вестник
  науки. 2025. № 12 (93). С. 256-275. URL:
  \url{https://cyberleninka.ru/article/n/sravnitelnyy-analiz-sovremennyh-sistem-avtomatizirovannogo-upravleniya-mikroklimatom-v-teplitse.}
\item
  Халина Т. М., Цыбанюк Д. Е. Разработка автоматизированной системы
  управления микроклиматом тепличного комплекса // Наука и молодежь:
  материалы XVII Всероссийской научно-технической конференции. Барнаул:
  АлтГТУ, 2020. С. 125-126.
\item
  Тарков Ю. М., Сукьясов С. В. Построение модели управления параметрами
  микроклимата в теплице // Материалы Иркутского ГАУ. 2025. С. 836-839.
\item
  Воробьева Н. С. {[}и др.{]}. Интеллектуальная система контроля
  температуры в теплице // Известия Нижневолжского агроуниверситетского
  комплекса. 2026. № 2 (86). С. 517-522.
\item
  Кельсин А. А. Автоматизированная система управления тепличным
  хозяйством на базе микроконтроллеров: выпускная квалификационная
  работа. Екатеринбург: Уральский государственный педагогический
  университет, 2021.
\item
  Гречиха А. С. Обзор системы управления микроклиматом
  автоматизированной теплицы для выращивания микрозелени // Агротехника
  и энергообеспечение. 2023. № 2.
\item
  Alhnaity B., Pearson S., Leontidis G., Kollias S. Using Deep Learning
  to Predict Plant Growth and Yield in Greenhouse Environments.
  arXiv:1907.00624. 2019. URL: \url{https://arxiv.org/abs/1907.00624.}
\item
  Gong L., Yu M., Jiang S., Cutsuridis V., Pearson S. Deep Learning
  Based Prediction on Greenhouse Crop Yield Combined TCN and RNN //
  Sensors. 2021. Vol. 21, No.~13. Article 4537. DOI: 10.3390/s21134537.
\item
  Ahn J. Y., Kim Y., Park H., Park S. H., Suh H. K. Evaluating
  Time-Series Prediction of Temperature, Relative Humidity, and CO\(_2\)
  in the Greenhouse with Transformer-Based and RNN-Based Models //
  Agronomy. 2024. Vol. 14, No.~3. Article 417. DOI:
  10.3390/agronomy14030417.
\item
  Huang S., Liu Q., Wu Y., Chen M., Yin H., Zhao J. Edible Mushroom
  Greenhouse Environment Prediction Model Based on Attention CNN-LSTM //
  Agronomy. 2024. Vol. 14, No.~3. Article 473. DOI:
  10.3390/agronomy14030473.
\item
  Choi W.-J., Yang M. Probabilistic Deep Learning Framework for
  Greenhouse Microclimate Prediction with Time-Varying Uncertainty and
  Covariance Analysis // Agriculture. 2025. Vol. 15, No.~23. Article
  2461. DOI: 10.3390/agriculture15232461.
\item
  Aborujilah A., Al-Sarem M., Abu-Zanona M. A. Forecast-Driven Climate
  Control for Smart Greenhouses: Energy Optimization Using LSTM Model //
  Energies. 2025. Vol. 18, No.~21. Article 5821. DOI:
  10.3390/en18215821.
\item
  Soumo E. A., Calisti M., Polydoros A. Data-Driven Greenhouse Climate
  Regulation in Lettuce Cultivation Using BiLSTM and GRU Predictive
  Control. arXiv:2507.21669. 2026. URL:
  \url{https://arxiv.org/abs/2507.21669.}
\item
  Thwin K. M. M., Horanont T., Phatrapornnant T. Machine-Learning
  Microclimate Forecasting for Adaptive Equipment Control via Web
  Integration in Open-Ventilated Greenhouses // AgriEngineering. 2024.
  Vol. 6, No.~3. P. 2845-2869. DOI: 10.3390/agriengineering6030165.
\item
  Mallick S., Airaldi F., Dabiri A., Sun C., De Schutter B.
  Reinforcement learning-based model predictive control for greenhouse
  climate control // Smart Agricultural Technology. 2025. Vol. 10.
  Article 100751. DOI: 10.1016/j.atech.2024.100751.
\item
  Msaad S., Harraway M., McAllister R. D. RL-Guided MPC for Autonomous
  Greenhouse Control. arXiv:2506.13278. 2025. URL:
  \url{https://arxiv.org/abs/2506.13278.}
\item
  Morcego B., Yin W., Boersma S., van Henten E., Puig V., Sun C.
  Reinforcement Learning Versus Model Predictive Control on Greenhouse
  Climate Control. arXiv:2303.06110. 2023. URL:
  \url{https://arxiv.org/abs/2303.06110.}
\item
  Xiao M., Lan J., Yu J., Ma W., Xie Q., Sun C. Grower-in-the-Loop
  Interactive Reinforcement Learning for Greenhouse Climate Control.
  arXiv:2505.23355. 2025. URL: \url{https://arxiv.org/abs/2505.23355.}
\item
  Platero-Horcajadas M., Pardo-Pina S., Cámara-Zapata J.-M.,
  Brenes-Carranza J.-A., Ferrández-Pastor F.-J. Enhancing Greenhouse
  Efficiency: Integrating IoT and Reinforcement Learning for Optimized
  Climate Control // Sensors. 2024. Vol. 24, No.~24. Article 8109. DOI:
  10.3390/s24248109.
\item
  Jawad M., Wahid F., Ali S., Ma Y., Alkhyyat A., Khan J., Lee Y. Energy
  optimization and plant comfort management in smart greenhouses using
  the artificial bee colony algorithm // Scientific Reports. 2025. DOI:
  10.1038/s41598-024-84141-5.
\item
  Teja G. S., Sathish P. MQTT Protocol based Smart Greenhouse
  Environment Monitoring System using Machine Learning // International
  Journal of Innovative Technology and Exploring Engineering. 2020. Vol.
  9, No.~9. P. 278-283. DOI: 10.35940/ijitee.I7149.079920.
\item
  Joni I. M. {[}et al.{]}. Smart Farming Experiment: IoT-Enhanced
  Greenhouse Design for Rice Cultivation with Foliar and Soil
  Fertilization // AgriEngineering. 2025. Vol. 7, No.~11. Article 380.
  DOI: 10.3390/agriengineering7110380.
\item
  Wu Y., Fu X. Clean and Smart Energy Technologies for Agricultural
  Energy Internet Systems: A Comprehensive Review and Future
  Perspectives // Applied Sciences. 2026. Vol. 16, No.~10. Article 4859.
  DOI: 10.3390/app16104859.
\item
  Hindi I., Alsharkawi A., Al-Ajlouni M., Qarallah B. Enhancing
  autonomous agriculture control systems in greenhouses for sustainable
  resource usage using deep learning techniques // PLOS One. 2026. DOI:
  10.1371/journal.pone.0344946.
\item
  Тюхтий Ю. А. Разработка интеллектуальной системы управления
  температурой в помещении: магистерская диссертация. Екатеринбург:
  Уральский федеральный университет, 2021.
\item
  Chen S., Liu A., Tang F., Hou P., Lu Y., Yuan P. A Review of
  Environmental Control Strategies and Models for Modern Agricultural
  Greenhouses // Sensors. 2025. Vol. 25, No.~5. Article 1388. DOI:
  10.3390/s25051388.
\item
  Ogunlowo Q. O., Akpenpuun T. D., Na W.-H., Rabiu A., Adesanya M. A.,
  Addae K. S., Kim H.-T., Lee H.-W. Analysis of Heat and Mass
  Distribution in a Single- and Multi-Span Greenhouse Microclimate //
  Agriculture. 2021. Vol. 11, No.~9. Article 891. DOI:
  10.3390/agriculture11090891.
\item
  Li H., Li A., Hou Y., Zhang C., Guo J., Li J., Ma Y., Wang T., Yin Y.
  Analysis of Heat and Humidity in Single-Slope Greenhouses with Natural
  Ventilation // Buildings. 2023. Vol. 13, No.~3. Article 606. DOI:
  10.3390/buildings13030606.
\item
  Ghaderi M., Reddick C., Sorin M. A Systematic Heat Recovery Approach
  for Designing Integrated Heating, Cooling, and Ventilation Systems for
  Greenhouses // Energies. 2023. Vol. 16, No.~14. Article 5493. DOI:
  10.3390/en16145493.
\item
  Самарин Г. Н., Попов А. Н., Плотников С. В., Ружьев В. А. Методика
  построения имитационной модели оптимального управления параметрами
  микроклимата сельскохозяйственного объекта // Аграрный научный журнал.
  2025. № 1. С. 126-135. DOI: 10.28983/asj.y2025i1pp126-135.
\end{enumerate}


\end{document}
